\documentclass[a4paper,landscape,8pt]{extarticle}
\usepackage[landscape,margin=0.7cm,top=0.5cm,bottom=0.7cm]{geometry}
\usepackage[utf8]{inputenc}
\usepackage[T1]{fontenc}
\usepackage[ngerman]{babel}
\usepackage{helvet}
\renewcommand{\familydefault}{\sfdefault}
\usepackage{xcolor}
\usepackage{enumitem}
\usepackage{multicol}
\usepackage{array}
\usepackage{tabularx}
\usepackage{fancyhdr}
\usepackage{amssymb}
\usepackage{graphicx}

% Farben
\definecolor{appleblue}{HTML}{007AFF}
\definecolor{applegray}{HTML}{8E8E93}
\definecolor{lightgray}{HTML}{F5F5F5}
\definecolor{bordergray}{HTML}{CCCCCC}
\definecolor{warnred}{HTML}{C62828}

\pagestyle{fancy}
\fancyhf{}
\renewcommand{\headrulewidth}{0pt}
\fancyfoot[C]{\footnotesize Seite \thepage{} von 3}

\setlength{\parindent}{0pt}
\setlength{\parskip}{2pt}

% Kompakte Listen
\setlist{nosep, leftmargin=1em, itemsep=0pt, topsep=1pt}

\newcommand{\infobox}[1]{%
\begin{center}
\fcolorbox{appleblue}{appleblue!8}{%
\begin{minipage}{0.95\linewidth}
\vspace{2mm}
#1
\vspace{2mm}
\end{minipage}}
\end{center}}

\newcommand{\warnbox}[1]{%
\fcolorbox{warnred}{warnred!8}{%
\begin{minipage}{0.95\linewidth}
\vspace{2mm}
\textbf{Wichtig:} #1
\vspace{2mm}
\end{minipage}}}

\begin{document}

% ============================================================================
% SEITE 1: WARUM & WIE
% ============================================================================

\begin{center}
{\LARGE\bfseries Meine Geräte-Inventur}\\[4pt]
{\large Welche Geräte besitze ich? Wo sind sie angemeldet?}
\end{center}

\vspace{1mm}
\hrule
\vspace{1mm}

\begin{multicols}{2}

\section*{Warum ist das wichtig?}

\infobox{
Stell dir vor, du rufst beim Apple-Support an. Die erste Frage ist immer:\\[2pt]
\textbf{\glqq Welches Gerät haben Sie? Welches Modell?\grqq{}}\\[2pt]
Wenn du das nicht weißt, kann dir niemand helfen.
}

\textbf{Diese Inventur hilft dir:}
\begin{itemize}
\item Bei Support-Anfragen die richtigen Infos parat haben
\item Wissen, welche Apple-ID wo angemeldet ist
\item Überblick behalten: Was gehört zusammen?
\item Bei Verlust/Diebstahl: Seriennummer für Polizei
\item Beim Verkauf: Wissen, was du hast
\end{itemize}

\vspace{2mm}

\section*{Was du herausfinden wirst}

Für jedes Gerät notierst du:

\begin{tabular}{@{}p{4cm}p{8cm}@{}}
\textbf{Gerätetyp} & iPhone, iPad, Mac, Apple Watch... \\[2mm]
\textbf{Modellname} & z.B. \glqq iPhone 13\grqq{} oder \glqq MacBook Air\grqq{} \\[2mm]
\textbf{Speicher} & z.B. 64 GB, 128 GB, 256 GB \\[2mm]
\textbf{Seriennummer} & Eindeutige Nummer für Support/Polizei \\[2mm]
\textbf{Apple-ID} & Welches Konto ist angemeldet? \\[2mm]
\textbf{iOS/macOS Version} & Welche Software-Version läuft? \\[2mm]
\end{tabular}

\columnbreak

\section*{Wo finde ich diese Infos?}

\subsection*{Auf dem iPhone / iPad:}

\textbf{Einstellungen} → \textbf{Allgemein} → \textbf{Info}

Dort findest du:
\begin{itemize}
\item Modellname (z.B. \glqq iPhone 13\grqq{})
\item Softwareversion (z.B. \glqq iOS 17.2\grqq{})
\item Seriennummer
\item Kapazität (Speicher)
\end{itemize}

\vspace{2mm}

\textbf{Apple-ID finden:}\\
\textbf{Einstellungen} → ganz oben steht dein Name → darunter die E-Mail-Adresse = deine Apple-ID

\vspace{4mm}

\subsection*{Auf dem Mac:}

\textbf{Apple-Menü} (Apfel-Symbol oben links) → \textbf{Über diesen Mac}

Dort findest du:
\begin{itemize}
\item Modellname (z.B. \glqq MacBook Air\grqq{})
\item Chip/Prozessor
\item Arbeitsspeicher
\item Seriennummer (auf \glqq Weitere Infos\grqq{} klicken)
\end{itemize}

\vspace{2mm}

\textbf{Apple-ID finden:}\\
\textbf{Systemeinstellungen} → \textbf{Apple-ID} (oben)

\end{multicols}

\vspace{1mm}

\begin{center}
\warnbox{Bewahre diese ausgefüllte Inventur an einem sicheren Ort auf – z.B. in einem Ordner zu Hause. NICHT auf dem Gerät selbst speichern!}
\end{center}

\newpage

% ============================================================================
% SEITE 2: INVENTAR-FORMULARE
% ============================================================================

\begin{center}
{\LARGE\bfseries Mein Geräte-Inventar}\\[4pt]
{\large Fülle für jedes Gerät einen Abschnitt aus}
\end{center}

\vspace{1mm}
\hrule
\vspace{2mm}

% Gerät 1
\begin{minipage}[t]{0.48\linewidth}
\fcolorbox{appleblue}{white}{%
\begin{minipage}{\linewidth}
\vspace{3mm}
\begin{center}
{\large\bfseries\textcolor{appleblue}{Gerät 1: iPhone}}
\end{center}
\vspace{2mm}

\renewcommand{\arraystretch}{1.5}
\begin{tabular}{@{}p{4cm}p{6.5cm}@{}}
\textbf{Modellname:} & \dotfill \\
\textbf{Farbe:} & \dotfill \\
\textbf{Speicher (GB):} & \dotfill \\
\textbf{iOS-Version:} & \dotfill \\
\textbf{Seriennummer:} & \dotfill \\
\textbf{Telefonnummer:} & \dotfill \\
\textbf{Apple-ID (E-Mail):} & \dotfill \\
\textbf{Mobilfunkanbieter:} & \dotfill \\
\end{tabular}
\renewcommand{\arraystretch}{1}
\vspace{3mm}
\end{minipage}}
\end{minipage}%
\hfill%
% Gerät 2
\begin{minipage}[t]{0.48\linewidth}
\fcolorbox{appleblue}{white}{%
\begin{minipage}{\linewidth}
\vspace{3mm}
\begin{center}
{\large\bfseries\textcolor{appleblue}{Gerät 2: iPad}}
\end{center}
\vspace{2mm}

\renewcommand{\arraystretch}{1.5}
\begin{tabular}{@{}p{4cm}p{6.5cm}@{}}
\textbf{Modellname:} & \dotfill \\
\textbf{Farbe:} & \dotfill \\
\textbf{Speicher (GB):} & \dotfill \\
\textbf{iPadOS-Version:} & \dotfill \\
\textbf{Seriennummer:} & \dotfill \\
\textbf{Mobilfunk (ja/nein):} & \dotfill \\
\textbf{Apple-ID (E-Mail):} & \dotfill \\
\textbf{Apple Pencil (ja/nein):} & \dotfill \\
\end{tabular}
\renewcommand{\arraystretch}{1}
\vspace{3mm}
\end{minipage}}
\end{minipage}

\vspace{3mm}

% Gerät 3
\begin{minipage}[t]{0.48\linewidth}
\fcolorbox{appleblue}{white}{%
\begin{minipage}{\linewidth}
\vspace{3mm}
\begin{center}
{\large\bfseries\textcolor{appleblue}{Gerät 3: Mac (Computer)}}
\end{center}
\vspace{2mm}

\renewcommand{\arraystretch}{1.5}
\begin{tabular}{@{}p{4cm}p{6.5cm}@{}}
\textbf{Modellname:} & \dotfill \\
\textbf{Baujahr:} & \dotfill \\
\textbf{Speicher (GB/TB):} & \dotfill \\
\textbf{macOS-Version:} & \dotfill \\
\textbf{Seriennummer:} & \dotfill \\
\textbf{Chip/Prozessor:} & \dotfill \\
\textbf{Apple-ID (E-Mail):} & \dotfill \\
\textbf{Benutzername:} & \dotfill \\
\end{tabular}
\renewcommand{\arraystretch}{1}
\vspace{3mm}
\end{minipage}}
\end{minipage}%
\hfill%
% Gerät 4
\begin{minipage}[t]{0.48\linewidth}
\fcolorbox{appleblue}{white}{%
\begin{minipage}{\linewidth}
\vspace{3mm}
\begin{center}
{\large\bfseries\textcolor{appleblue}{Gerät 4: Apple Watch}}
\end{center}
\vspace{2mm}

\renewcommand{\arraystretch}{1.5}
\begin{tabular}{@{}p{4cm}p{6.5cm}@{}}
\textbf{Modellname:} & \dotfill \\
\textbf{Größe (mm):} & \dotfill \\
\textbf{Farbe/Material:} & \dotfill \\
\textbf{watchOS-Version:} & \dotfill \\
\textbf{Seriennummer:} & \dotfill \\
\textbf{Mobilfunk (ja/nein):} & \dotfill \\
\textbf{Gekoppelt mit iPhone:} & \dotfill \\
\textbf{Apple-ID (E-Mail):} & \dotfill \\
\end{tabular}
\renewcommand{\arraystretch}{1}
\vspace{3mm}
\end{minipage}}
\end{minipage}

\vspace{3mm}

% Zusätzliches Gerät
\begin{minipage}[t]{0.48\linewidth}
\fcolorbox{applegray}{white}{%
\begin{minipage}{\linewidth}
\vspace{3mm}
\begin{center}
{\large\bfseries\textcolor{applegray}{Weiteres Gerät}}
\end{center}
\vspace{2mm}

\renewcommand{\arraystretch}{1.5}
\begin{tabular}{@{}p{4cm}p{6.5cm}@{}}
\textbf{Gerätetyp:} & \dotfill \\
\textbf{Modellname:} & \dotfill \\
\textbf{Speicher:} & \dotfill \\
\textbf{Software-Version:} & \dotfill \\
\textbf{Seriennummer:} & \dotfill \\
\textbf{Apple-ID (E-Mail):} & \dotfill \\
\textbf{Notizen:} & \dotfill \\
& \dotfill \\
\end{tabular}
\renewcommand{\arraystretch}{1}
\vspace{3mm}
\end{minipage}}
\end{minipage}%
\hfill%
% Meine Apple-ID Übersicht
\begin{minipage}[t]{0.48\linewidth}
\fcolorbox{warnred}{warnred!5}{%
\begin{minipage}{\linewidth}
\vspace{3mm}
\begin{center}
{\large\bfseries\textcolor{warnred}{Meine Apple-ID(s)}}
\end{center}
\vspace{2mm}

\renewcommand{\arraystretch}{1.5}
\begin{tabular}{@{}p{4cm}p{6.5cm}@{}}
\textbf{Haupt-Apple-ID:} & \dotfill \\
\textbf{Passwort-Hinweis:} & \dotfill \\
\textbf{Vertrauensperson:} & \dotfill \\
\textbf{Weitere Apple-ID:} & \dotfill \\
\textbf{(falls vorhanden)} & \\
\end{tabular}
\renewcommand{\arraystretch}{1}

\vspace{2mm}
\footnotesize\textit{Passwort NIEMALS hier aufschreiben! Nur einen Hinweis, der dir hilft, dich zu erinnern.}
\vspace{3mm}
\end{minipage}}
\end{minipage}

\newpage

% ============================================================================
% SEITE 3: CHECKLISTE & TIPPS
% ============================================================================

\begin{center}
{\LARGE\bfseries Checkliste \& Tipps}\\[4pt]
{\large Hast du alles gefunden?}
\end{center}

\vspace{2mm}
\hrule
\vspace{6mm}

\begin{multicols}{2}

\section*{Checkliste: Das habe ich herausgefunden}

\begin{itemize}[label=$\square$, itemsep=6pt]
\item Ich weiß, welches \textbf{iPhone-Modell} ich habe
\item Ich kenne die \textbf{iOS-Version} meines iPhones
\item Ich habe die \textbf{Seriennummer} meines iPhones notiert
\item Ich weiß, welche \textbf{Apple-ID} auf meinem iPhone angemeldet ist
\item Ich habe meine weiteren Geräte (iPad, Mac...) dokumentiert
\item Ich weiß, ob auf allen Geräten die \textbf{gleiche Apple-ID} angemeldet ist
\item Ich habe diese Inventur an einem \textbf{sicheren Ort} aufbewahrt
\end{itemize}

\vspace{6mm}

\section*{Häufige Fragen}

\textbf{Warum haben manche Geräte verschiedene Apple-IDs?}\\
Das passiert manchmal, wenn z.B. ein Gerät gebraucht gekauft wurde oder jemand anderes es eingerichtet hat. Idealerweise sollte auf allen DEINEN Geräten die GLEICHE Apple-ID sein.

\vspace{4mm}

\textbf{Was ist der Unterschied zwischen Modell und Seriennummer?}\\
Das \textbf{Modell} (z.B. \glqq iPhone 13\grqq{}) beschreibt den Typ – davon gibt es Millionen gleiche.\\
Die \textbf{Seriennummer} ist einzigartig – wie ein Fingerabdruck für dein Gerät.

\columnbreak

\section*{Wo finde ich die Info nochmal?}

\begin{tabular}{@{}p{3cm}p{7cm}@{}}
\textbf{iPhone/iPad:} & Einstellungen → Allgemein → Info \\[3mm]
\textbf{Apple-ID:} & Einstellungen → [Dein Name] oben \\[3mm]
\textbf{Mac:} & Apfel-Menü → Über diesen Mac \\[3mm]
\textbf{Apple Watch:} & Watch-App auf iPhone → Allgemein → Info \\[3mm]
\end{tabular}

\vspace{6mm}

\section*{Im Notfall: Gerät verloren?}

Wenn du diese Inventur ausgefüllt hast, kannst du:

\begin{enumerate}[itemsep=3pt]
\item Die \textbf{Seriennummer} der Polizei geben
\item Über \textbf{iCloud.com} dein Gerät orten (\glqq Wo ist?\grqq{})
\item Beim Apple-Support mit den richtigen Infos anrufen
\end{enumerate}

\vspace{4mm}

\infobox{
\textbf{Tipp:} Mache ein Foto von dieser ausgefüllten Seite und speichere es in deiner E-Mail an dich selbst – so hast du die Infos auch digital.
}

\end{multicols}

\vspace{4mm}

\begin{center}
\fcolorbox{bordergray}{lightgray}{%
\begin{minipage}{0.95\linewidth}
\centering
\vspace{3mm}
\textbf{Datum der Inventur:} \rule{4cm}{0.4pt} \hspace{2cm} \textbf{Nächste Aktualisierung:} \rule{4cm}{0.4pt}
\vspace{3mm}
\end{minipage}}
\end{center}

\vfill

\begin{center}
\footnotesize\textit{Stand: \today{} – Aktualisiere diese Liste, wenn du ein neues Gerät bekommst oder eines verkaufst.}
\end{center}

\end{document}
